\documentclass[12pt,a4paper]{article}
\usepackage[utf8]{inputenc}
\usepackage[T1]{fontenc}
\usepackage[brazil]{babel}
\usepackage{amsmath,amssymb}
\usepackage{geometry}
\geometry{margin=2.5cm}
\title{Material de Estudo: Apresenta\c{c}\~ao de Trabalhos III --- Avalia\c{c}\~ao por Pares}
\author{Princ\'ipio de Modelagem Matem\'atica}
\date{Unidade 34}
\begin{document}
\maketitle

\section*{Objetivo}
Desenvolver o senso cr\'itico cient\'ifico por meio da avalia\c{c}\~ao construtiva das apresenta\c{c}\~oes dos colegas.

\section*{Por qu\^e Avaliar os Colegas?}

\begin{itemize}
  \item Exercita a leitura cr\'itica --- habilidade central na ci\^encia.
  \item Ajuda quem apresenta a melhorar o trabalho.
  \item Faz quem avalia pensar mais profundamente sobre os m\'etodos.
  \item Prepara para o processo de revis\~ao por pares em peri\'odicos.
\end{itemize}

\section*{Formul\'ario de Avalia\c{c}\~ao por Pares}

Para cada apresenta\c{c}\~ao assistida, preencha:

\subsection*{Informa\c{c}\~oes Gerais}
\begin{itemize}
  \item Grupo avaliado: \underline{\hspace{6cm}}
  \item Data: \underline{\hspace{4cm}}
  \item Avaliador: \underline{\hspace{5cm}}
\end{itemize}

\subsection*{Avalia\c{c}\~ao (1 = ruim; 5 = excelente)}

\begin{tabular}{|l|c|}
\hline
\textbf{Crit\'erio} & \textbf{Nota (1--5)} \\
\hline
Clareza da pergunta de pesquisa & \\
Adequa\c{c}\~ao do modelo ao problema & \\
Explica\c{c}\~ao dos m\'etodos e hip\'oteses & \\
Qualidade das figuras e resultados & \\
Profundidade da discuss\~ao & \\
Clareza das conclus\~oes & \\
Comunica\c{c}\~ao oral e visual & \\
Respostas \`as perguntas & \\
\hline
\textbf{M\'edia} & \\
\hline
\end{tabular}

\subsection*{Perguntas de Reflex\~ao}
Responda em 2--4 frases cada:
\begin{enumerate}
  \item Qual foi o ponto mais forte da apresenta\c{c}\~ao?
  \item Qual foi a limita\c{c}\~ao mais importante n\~ao discutida pelo grupo?
  \item O que voc\^e faria diferente no modelo ou na implementa\c{c}\~ao?
  \item A valida\c{c}\~ao apresentada foi convincente? Por qu\^e?
\end{enumerate}

\section*{Exerc\'icios}

\begin{enumerate}
  \item \textbf{(Feedback)} Para a apresenta\c{c}\~ao que voc\^e assistiu: escreva um par\'agrafo de feedback construtivo (o que foi bom e o que poderia melhorar, com exemplos concretos).
  
  \item \textbf{(Quest\~ao} cient\'ifica) Formule UMA pergunta cient\'ifica relevante sobre o modelo apresentado que o grupo n\~ao respondeu.
  
  \item \textbf{(Compara\c{c}\~ao)} Compare o modelo do grupo avaliado com um modelo similar que vimos na disciplina. Quais as semelhan\c{c}as e diferen\c{c}as metodol\'ogicas?
  
  \item \textbf{(Extens\~ao)} Sugira uma extens\~ao do trabalho apresentado que seria realiz\'avel e interessante. Justifique.
  
  \item \textbf{(Reflex\~ao)} Depois de ouvir todas as apresenta\c{c}\~oes: qual abordagem de modelagem voc\^e achou mais elegante ou eficaz? Por qu\^e?
  
  \item \textbf{(Autofeedback)} Depois de ouvir os colegas: o que voc\^e mudaria na sua pr\'opria apresenta\c{c}\~ao?
  
  \item \textbf{(Debate)} Escreva um argumento a favor e um argumento contra a seguinte afirma\c{c}\~ao: ``Modelos simples s\~ao sempre prefer\'iveis a modelos complexos.''
\end{enumerate}

\section*{Checklist de Participa\c{c}\~ao}
\begin{itemize}
  \item[$\square$] Assistiu \`as apresenta\c{c}\~oes com aten\c{c}\~ao.
  \item[$\square$] Fez ao menos uma pergunta durante as apresenta\c{c}\~oes.
  \item[$\square$] Preencheu o formul\'ario de avalia\c{c}\~ao com cuidado.
  \item[$\square$] Feedback \'e construtivo, n\~ao apenas cr\'itico.
  \item[$\square$] Perguntas s\~ao cient\'ificas, n\~ao s\'o formais.
\end{itemize}

\section*{Resumo R\'apido}
\begin{itemize}
  \item Avalia\c{c}\~ao por pares: treinamento cr\'itico e colaborativo.
  \item Feedback construtivo: aponte o que funciona e sugira melhorias concretas.
  \item Perguntas cient\'ificas s\~ao mais valiosas que perguntas formais.
\end{itemize}

\section*{Refer\^encias}
\begin{itemize}
  \item Notas de aula da disciplina.
\end{itemize}

\end{document}
